\chapter*{Core Operating System Overview}
\addcontentsline{toc}{chapter}{Core Operating System Overview}

\section{Introduction}
This section gives a high-level view of the core operating system we built across
milestones M1 through M6.  The system runs on ARMv8 (AArch64) hardware either
emulated on QEMU or through a 4-core Toradex iMX8X board and is based off the Barrelfish
OS. The detailed design and implementation of each subsystem are described in
the milestone chapters that follow. This overview is a brief, high-level overview 
of the different components we implemented and how they interact with the existing
Barrelfish structure.

% TODO
% \begin{figure}[H]
%     \centering
%     \includegraphics[width=0.5\textwidth]{}
%     \caption{Core OS Architecture}
%     \label{fig:OSarch}
% \end{figure}


\section{Microkernel and Capabilities}

\subsection{Minimal Trusted Core}
The Barrelfish kernel, called the CPU driver, is a thin privileged layer that runs 
independently on each core. Each core's CPU driver manages only the hardware directly 
beneath it(the MMU, the interrupt controller, and the scheduling of user-space dispatchers).
Operations that are supposed to be executed by the kernel in usual monolithic kernels
are now handled by user-space programs. As such the kernel in Barrelfish is a microkernel.
Microkernels only handle essential core operations and push functionality out of the
kernel. 

\medskip\noindent
When a user-space program needs to perform a privileged action, say map a page,
create an endpoint, boot a second core, it issues a system call that
names a specific capability it already holds.  The kernel checks that the
capability grants the requested permission and, if so, carries out the operation
atomically. These type of capabilities are called partitioned capabilities.

\subsection{Capabilities}

Capabilities are essentially tokens that are associated with kernel operations/objects
and represent the rights held over them. There are different types of capabilties that
are retyped from physical memory. Namely, the following are commonly utilized within
our implementation.

\begin{itemize}[nosep]
  \item \textbf{RAM} --- a contiguous range of physical memory that has not yet
        been committed to any purpose.  A process that holds a RAM capability
        can retype it into frames, page-table nodes, endpoints, or other objects.
  \item \textbf{Frame} --- a mappable physical page. Can be installed in a page
        table or shared with another process by transferring the capability.
  \item \textbf{VNode} --- a hardware page-table node (L0--L3 on ARMv8). Holds
        the mappings of a particular level of the address-space tree.
  \item \textbf{Endpoint} --- a communication buffer in the kernel. The target
        of a Local Message-Passing (LMP) send.
  \item \textbf{Dispatcher} --- the scheduler context of a running domain.
        Holding this capability allows pausing, resuming, or destroying the
        domain.
  \item \textbf{Kernel Control Block (KCB)} --- per-core kernel state. Must be
        allocated and handed to the kernel before a new core can be started.
\end{itemize}

\medskip\noindent
Capabilities live in per-process CSpaces, capability namespaces
structured as trees of CNodes, and are never directly accessible to user code.
To allocate capabilities and retype them into the correct capabilities, we acquire
chunk of untyped memory, split the memory into the correct size, and then retype the
capability into the desired type.


\section{User-Space System Structure}


\subsection{The \co{init} Process}

On each core a single privileged user-space process called \co{init} acts
as the system manager.  It is the first program to run after the CPU driver
hands off execution, and it holds the initial RAM capabilities granted by 
the bootloader.

\medskip\noindent
\co{init} is responsible for:

\begin{itemize}[nosep]
  \item Dividing physical memory among itself and the processes it spawns.
  \item Constructing the virtual address space of each new process and loading
        its ELF image.
  \item Providing every OS service that other processes request via RPC.
        Memory allocation, serial I/O, process lifecycle operations, and
        coordination with peer cores.
\end{itemize}

\medskip\noindent
Because \co{init} holds the only RAM capabilities in the system, no process
can obtain memory without going through it. This gives \co{init} complete
control over resources without any kernel support.

\subsection{Child Processes}

All other programs run as children of \co{init}.  A newly spawned process
starts with almost nothing: a valid address space, a stack, its ELF image
loaded, and the endpoint of \co{init}'s RPC server(which is a capability).
Everything else the process needs (more memory, I/O access, knowledge of other
processes) is obtained by sending requests to \co{init} over that channel.

\medskip\noindent
Spawning, pausing, resuming, and killing processes are all \co{init}
responsibilities, carried out through the RPC interface(described later).  

\section{Memory Management}


Physical and virtual memory management are both implemented entirely in user
space.

\subsection{Physical Memory}

\co{init} holds a pool of RAM capabilities passed up from the bootloader.
When a process requests memory, \co{init}'s physical allocator carves a
sub-region out of the pool, retypes it into a Frame capability, and transfers
that capability over LMP to the requesting process.  The requesting process then
maps the frame into its own address space using the VNode capabilities it holds.

\medskip\noindent
The kernel is involved only at the retyping step (to enforce the capability type
system) and at the mapping step (to update the hardware MMU). Every other desicion is
made entirely in user space through our code(size to allocate, which region to allocate, etc.).

\subsection{Virtual Memory}

Each process manages its own virtual address space.  The paging library
(\co{lib/aos/paging}) tracks which virtual address ranges are free and which
are occupied, allocates and maps page-table VNodes on demand, and installs a
page fault handler that lazily maps frames as they are first accessed.

\section{Inter-Process Communication}

In a microkernel, since kernel functionality is pushed out to user-space programs.
This means other programs can no longer invoke these services through a 
simple function call or a syscall that touches shared kernel state. Rather it must
send a message to the process that owns the service and wait for a reply. 
IPC is therefore is the fundamental mechanism through which the system functions
and is required by other processes.

\medskip\noindent
The system provides two IPC mechanisms, chosen based on whether the
communicating parties are on the same core or on different cores.

\subsection{Local Message Passing (LMP) --- Intra-Core}

LMP handles communication between processes on the same core. Messages are
deposited into a kernel-managed endpoint buffer and the receiving process is
scheduled when a message arrives. The entire RPC protocol between child
processes and \co{init} is built on top of LMP.

\subsection{User Message Passing (UMP) --- Inter-Core}

LMP endpoints cannot be shared across cores, so cross-core communication is
implemented in user space over a shared physical frame. This makes UMP
higher-latency than LMP, but it requires no kernel support and scales as cores
are added.

\section{The RPC Protocol}

Since \co{init} owns all system resources, every process that needs memory,
wants to spawn a child, or reads from serial I/O has to ask \co{init} for it.
RPC gives these interactions a clean structure. Rather than ad-hoc messages,
each service has a well-defined call and reply, making the interface between
processes predictable and easy to extend.

\medskip\noindent
All communication between a child process and \co{init} follows a unified
request-reply protocol, regardless of whether the underlying transport is LMP
or UMP. From a child process's perspective, making a service request looks the
same whether \co{init} is running on the same core or a different one. The RPC
library handles the details: it sends a tagged request, blocks until the
matching reply arrives, and returns the result to the caller. On the \co{init}
side, incoming messages are dispatched to the appropriate handler. If the
request concerns a process on a remote core, \co{init} forwards it over UMP and
relays the reply back transparently.

\section{Multi-Core}

The system runs on up to four cores, each hosting its own independent
\co{init} instance with its own memory pool and process table.  Cores
communicate over UMP channels backed by shared physical frames, forming a
fully-connected mesh. From the perspective of a user process, the
system appears as a single coherent OS regardless of how many cores are active.